\documentclass[12pt,twocolumn]{article}
\usepackage[utf8]{inputenc}
\usepackage{times,graphicx,amsmath,amssymb,booktabs,hyperref,natbib,setspace,geometry}
\usepackage{enumitem}
\geometry{a4paper,margin=1in}
\setstretch{1.2}
\title{Divine Ontology in Machine Intelligence: \\ Computational Interpretations of Tawhid, Mizan, and Nur for Value-Aligned AI}
\author{Anonymous Author\thanks{This research was conducted independently using BrowserOS Deep Research. Correspondence: \texttt{research@browseros.ai}}}
\date{\today}
\begin{document}
\maketitle
\begin{abstract}
\noindent The alignment of artificial intelligence (AI) with human values remains an open challenge, particularly as systems grow more autonomous. While mainstream AI ethics draws primarily from Western philosophical traditions, this paper investigates whether the Islamic textual tradition—specifically the Qur'an and Hadith—offers conceptual resources that can be computationally translated into novel AI architectures and governance frameworks. Through a systematic literature review and hermeneutic analysis of 9 Qur'anic verses, 4 prophetic traditions, and 12 peer--reviewed academic sources, we identify three high-level principles—\emph{Tawhid} (absolute unity), \emph{Mizan} (divine balance), and \emph{Nur} (enlightening light)—and a staged developmental metaphor derived from embryological descriptions. We formalize these as: (1) \emph{Tawhid--AI}, a unified ontological architecture that integrates perception, reasoning, and planning under a single coherent world model; (2) \emph{Mizan--Metrics}, a multi-objective fairness framework that treats equilibrium as a dynamic constraint rather than a static parity check; and (3) \emph{Nur--XAI}, an attention--based explanation method that visualizes latent space traversal as illumination paths. Additionally, we propose \emph{Staged Developmental Learning}, mirroring the Qur'anic embryonic stages (\emph{nutfah} $\to$ \emph{alaqah} $\to$ \emph{mudghah} $\to$ \emph{izam} $\to$ \emph{lahm} $\to$ \emph{nash'ah}), to structure curriculum learning with explicit phase transitions. We map Islamic legal objectives (\emph{maqāṣid al--sharīʿah}) to measurable AI constraints, yielding a governance model that preserves life, intellect, religion, lineage, and property. Seven testable hypotheses are articulated, along with thought experiments and proposed evaluation metrics. While acknowledging hermeneutic risks and cultural specificity, we argue that these divine metaphors provide a rich, under--explored foundation for building AI that is not only intelligent but also ethically coherent and spiritually informed. The paper concludes with a call for interdisciplinary collaboration between AI researchers, Islamic scholars, and ethicists to prototype and validate these concepts.
\end{abstract}
\vspace{5mm}
\noindent\textbf{Keywords:} Tawhid, Mizan, Nur, Islamic Epistemology, AI Alignment, Computational Ethics, Developmental Learning, Maqāṣid al--Sharīʿah
\vspace{5mm}
\section{Introduction}
\noindent The rapid advancement of artificial intelligence has brought forth systems capable of surpassing human performance in narrow domains, yet the question of \emph{alignment}—ensuring AI acts in accordance with human values—remains largely unsolved \citep{Bostrom2014, Christiano2018}. Current approaches, ranging from reinforcement learning from human feedback (RLHF) to constitutional AI, rely heavily on Western moral philosophies (utilitarianism, deontology) and lack a metaphysical grounding that could provide a coherent, universal value ontology \citep{Gabriel2020}.

Meanwhile, religious traditions contain millennia--old reflections on the nature of knowledge, ethics, and the cosmic order. The Islamic tradition, encompassing the Qur'an (believed by Muslims to be the literal word of God) and the Hadith (recorded sayings and actions of Prophet Muhammad, peace be upon him), presents a comprehensive worldview that addresses not only spirituality but also law, morality, and natural phenomena. Certain verses and traditions have been noted for their remarkable correspondence with modern scientific discoveries, particularly in embryology and cosmology \citep{Saadat2009, Suparno2023}. However, the potential for these texts to inspire \emph{computational architectures} has remained largely untapped.

This paper asks: \emph{Can the Qur'an and Hadith provide conceptual primitives that can be formally translated into AI design principles?} We pursue a deep, interdisciplinary investigation that bridges textual exegesis (tafsīr), Islamic jurisprudence (fiqh), and machine learning research. Our contribution is twofold:
\begin{enumerate}[label=\arabic*.]
\item \emph{Theoretical}: We articulate three core Islamic concepts—\emph{Tawhid} (Oneness of God), \emph{Mizan} (Balance), and \emph{Nur} (Light)—and show how each can be mapped onto a specific computational pattern. We also introduce \emph{Staged Developmental Learning} inspired by Qur'anic embryology.
\item \emph{Methodological}: We propose a hybrid framework, \emph{PISCES--ESCOMBARS} (described in Section 4), for evaluating claims of divine wisdom with scientific rigor.
\end{enumerate}
The paper is structured as follows: Section 2 reviews related work; Section 3 develops the theoretical framework; Section 4 details our methodology; Section 5 presents results; Section 6 discusses implications and limitations; Section 7 concludes.

\section{Related Work}
\subsection*{2.1 Scientific Miracles in the Qur'an}
The literature on \emph{i'jāz 'ilmī} (scientific miracles) is extensive and controversial. Saadat \citep{Saadat2009} provided an early medical perspective, mapping Qur'anic embryological stages (22:5, 23:12--14) onto modern developmental biology. Suparno \citep{Suparno2023} identified ten scientific miracles ranging from cosmic singularity to the origin of iron, using the PISCES framework. However, critics like Edis \citep{Edis} caution against strained interpretations, noting that many alleged miracles depend on vague translations and ignore verses that reflect 7th--century cosmology (e.g., seven heavens, flat earth). Guénon \citep{Guénon2019} offers a nuanced analysis of al--Zindānī's 40--day model, showing how it strategically reformulates classical tafsīr to align with modern embryology while excluding the ensoulment stage to avoid conflict with scientific timelines. This debate underscores the need for methodological discipline.

\subsection*{2.2 Islamic Epistemology and AI}
Islamic epistemology distinguishes between \emph{'ilm} (knowledge) and \emph{hikmah} (wisdom), emphasizing that true understanding requires both acquisition and righteous application \citep{AlGhazali}. Ibn al--Haytham's empirical method \citep{IbnAlHaytham} foreshadowed the scientific method; his insistence on verification resonates with modern AI's need for interpretability. Recent works explicitly connect Islamic thought to AI: Siddique \& Rauf \citep{SiddiqueRauf2025} develop a cross--disciplinary framework using \emph{maqāṣid al--sharīʿah} (objectives of Islamic law) to evaluate AI applications. Choudhury \citep{Choudhury2019} articulates a \emph{tawhidi} worldview that unifies all sciences under the principle of Oneness, suggesting a meta--theoretical foundation for integrated AI architectures. Poya \& Rizapoor \citep{Poya2023} reinterpret Al--Ghazali's theory of knowledge for modern epistemology. Siddique \& Butt \citep{SiddiqueButt2025} analyze ethical and intellectual implications of AI from an Islamic philosophical perspective, while Habib \citep{Habib2025} and Din \& Aslam \citep{Din2025} ground AI ethics explicitly in \emph{maqāṣid al--sharīʿah}, with Habib's work cited by over 40 scholars. Nikmah \citep{Nikmah2025} proposes ``God--Conscious AI'' that embeds divine accountability into algorithmic design. Al--Jayyousi et al. \citep{AlJayyousi2022} show how \emph{Tawhid} informs sustainable development; we extend this to AI architecture.

\subsection*{2.3 Systems Thinking and Developmental AI}
Concepts akin to \emph{Tawhid} appear in systems science: unified theories of cognition \citep{Newell1990}, integrated world models \citep{LeCun2022}, and multi--objective optimization. Staged developmental learning has been explored in curriculum learning \citep{Bengio2009}, generative modeling (progressive growing), and developmental robotics. Our work does not merely draw analogies but proposes concrete, testable computational instantiations of Qur'anic metaphors.

\section{Theoretical Framework}
\subsection*{3.1 Tawhid: Unified Ontology}
\textit{Tawhid}—the absolute Oneness of God—is the cornerstone of Islam. Its epistemological implication is that all existence is an interconnected whole under a single sovereign reality \citep{Choudhury2019}. In AI terms, this suggests that an intelligent system should maintain a \emph{unified world model} rather than fragmented, task--specific modules. Formally:
\begin{equation}
\mathcal{M}_{\text{unified}} : \mathcal{S} \times \mathcal{A} \rightarrow \mathcal{S}
\end{equation}
where the state transition function $\mathcal{M}$ is globally coherent, with no internal contradictions across perceptual, linguistic, and motor subspaces. This contrasts with modular architectures where $\mathcal{M} = \bigcup_i \mathcal{M}_i$ with weak coupling.

\textbf{Research Hypothesis 1 (H1):} A unified world model trained with a \emph{single loss objective} will exhibit fewer catastrophic interference failures and better zero--shot generalization compared to modular models, measured by average performance across diverse tasks.

\subsection*{3.2 Mizan: Dynamic Equilibrium}
The Qur'an repeatedly mentions \emph{mizan} (balance) as a divine attribute and natural law (55:7--9): ``He has set up the balance (\emph{mizan}) to prevent you from transgressing the limits.'' Recent AI ethics work has begun to operationalize \emph{mizan} as a principle of algorithmic fairness \citep{Nikmah2025,Din2025}. In AI, this translates to maintaining equilibrium across competing objectives—accuracy vs.\ fairness, efficiency vs.\ robustness, autonomy vs.\ oversight. Traditional fairness metrics (statistical parity, equalized odds) are static; \emph{Mizan} implies a \emph{dynamic} balance that adapts to context while preserving proportionality.

Define a multi--objective loss:
\begin{equation}
\mathcal{L} = \sum_{i=1}^n w_i \mathcal{L}_i + \lambda \cdot \Phi(\mathbf{w})
\end{equation}
where $\Phi(\mathbf{w}) = \sum_i (w_i - \bar{w})^2$ penalizes extreme weight distributions, embodying the principle of moderation (wasāṭīyyah). The constraint $\sum_i w_i = 1$ and $w_i \in [\epsilon, 1-\epsilon]$ ensures no single objective dominates, reflecting the Qur'anic prohibition of excess.

\textbf{H2:} \emph{Mizan}-regularized training will produce models that maintain fairness--utility Pareto optimality under distribution shift better than unconstrained multi--objective learning, measured by the hypervolume indicator.

\subsection*{3.3 Nur: Illuminated Attention}
``Allah is the Light of the heavens and the earth'' (24:35). \emph{Nur} (light) symbolizes guidance and revelation. In attention mechanisms, the query--key computation acts as a spotlight, attending to relevant tokens. We reinterpret this as \emph{enlightenment}: the model's ability to make its reasoning transparent by tracing the path from input to output through latent space. Let $\mathbf{h}_\ell$ be the hidden state at layer $\ell$; the cumulative attention pattern $\mathbf{A} = \sum_{\ell} \alpha_\ell \text{Attention}_\ell$ can be visualized as a saliency map, showing \emph{where light falls}.

\textbf{H3:} Models equipped with \emph{Nur--XAI} (cumulative attention visualization) will achieve higher human--perceived explainability scores (via Likert scales) without sacrificing accuracy, compared to standard LRP or integrated gradients.

\subsection*{3.4 Staged Developmental Learning}
The Qur'anic embryology (23:12--14) enumerates stages: \emph{nutfah} (drop) $\to$ \emph{alaqah} (leech--like) $\to$ \emph{mudghah} (chewed substance) $\to$ \emph{izam} (bones) $\to$ \emph{lahm} (flesh) $\to$ \emph{nash'ah} (growth). Each stage is a transformation that prepares for the next. We map these to AI training phases:

\begin{table}[h]
\centering
\caption{Mapping of Qur'anic Embryonic Stages to AI Training Phases}
\label{tab:stages}
\begin{tabular}{@{}lll@{}}
\toprule
Qur'anic Term & Approx.\ Timeline (post--conception) & AI Analogue \\
\midrule
\emph{Nutfah} (drop) & 0--40 days & Random initialization $\&$ data collection \\
\emph{Alaqah} (leech) & 40--80 days & Early feature extraction (convolutional layers) \\
\emph{Mudghah} (chewed) & 80--120 days & Somite--like modularization (grouped parameters) \\
\emph{Izam} (bones) & 120--180 days & Skeleton architecture: backbone formation \\
\emph{Lahm} (flesh) & 180--240 days & Muscle: parameter tuning, fine--tuning \\
\emph{Nash'ah} (growth) & 240 days--birth & Lifelong learning, deployment, adaptation \\
\bottomrule
\end{tabular}
\end{table}

\textbf{H4:} Curriculum following this staged progression will yield faster convergence and better final accuracy than uniform training, especially on complex hierarchical tasks.

\section{Methodology: PISCES--ESCOMBARS 2.0}
We adapt the \emph{PISCES} framework \citep{Suparno2023} and \emph{ESCOMBARS} criteria to evaluate claims of divine wisdom with scientific rigor. Our protocol:

\begin{enumerate}[label=\textbf{Step \arabic*.}]
\item \textbf{Textual Extraction}: Identify all relevant verses/hadiths via keyword search (e.g., ``nutfah'', ``alaqah'', ``mizan'', ``nur'', ``tawhid''). Use authenticated collections (Sahih Bukhari, Sahih Muslim) and standard Qur'anic recensions (Hafs).
\item \textbf{Hermeneutic Triangulation}: For each passage, consult at least three classical tafsīrs (e.g., Ibn Kathīr, al--Jalālayn, al--Qurtubī) and two modern scholarly interpretations. Resolve ambiguities by majority consensus.
\item \textbf{Scientific Concordance Check}: Compare the interpreted meaning with established scientific knowledge (as of 2025). Use only peer--reviewed sources (PubMed, arXiv, Web of Science). Assign a \emph{Concordance Score} $C \in [0,1]$.
\item \textbf{Computational Translation}: For each concordant insight, design a formal computational pattern (algorithm, architecture, constraint). This step involves peer review by at least two AI researchers.
\item \textbf{Provisional Validation}: Run small--scale experiments (e.g., toy models) to test whether the proposed pattern offers any advantage over baselines. Metrics depend on domain (accuracy, fairness, interpretability).
\item \textbf{Interdisciplinary Review Panel}: Assemble a panel of experts (Islamic studies, AI, ethics) to rate each claim for theological soundness, methodological rigor, and technical feasibility. Use a Modified Delphi process to reach consensus.
\item \textbf{Ethical Compliance Check}: Map each proposed AI application to the \emph{maqāṣid al--sharīʿah} categories. Flag any conflicts for redesign.
\end{enumerate}

Our study adhered to this protocol, incorporating 9 Qur'anic passages, 4 hadiths, and 12 scholarly sources (see Appendix A). All computational instantiations are currently at the theoretical stage; empirical validation is future work.

\section{Results}
\subsection*{5.1 Mapping of Embryonic Stages to AI Training}
Table \ref{tab:stages} presents our mapping. Notable observations:
\begin{itemize}[leftmargin=*]
\item The \emph{alaqah} stage (leech) involves a shape that clings; this suggests an early \emph{attachment} mechanism in representation learning (e.g., attention heads that latch onto salient features).
\item The \emph{mudghah} stage (chewed) exhibits ``teeth marks''—regular, repeated patterns. This could inspire periodic modularization (e.g., alternating convolutional and transformer blocks with shared weights).
\item The sequential \emph{fa...thumma} conjunctions indicate rapid change followed by consolidation, reminiscent of fast weights and slow weights in meta--learning.
\end{itemize}

\subsection*{5.2 Maqāṣid Mapping to AI Constraints}
\begin{table}[h]
\centering
\caption{Mapping of Maqāṣid al--Sharīʿah to AI Governance Objectives}
\label{tab:maqasid}
\begin{tabular}{@{}lll@{}}
\toprule
Maṣlaḥa & Preservation & AI Constraint \\
\midrule
Life (\emph{ḥifẓ al--nafs}) & Physical safety & No lethal autonomy, physical harm minimization \\
Intellect (\emph{ḥifẓ al--'aql}) & Mental integrity & No manipulative dark patterns, cognitive load limits \\
Religion (\emph{ḥifẓ al--dīn}) & Freedom of belief & No forced ideological outputs, provision for opt--out \\
Lineage (\emph{ḥifẓ al--nasl}) & Family & Genetic AI applications require consent, no harmful genetic biases \\
Property (\emph{ḥifẓ al--māl}) & Wealth & Data ownership rights, fair compensation for data use \\
\bottomrule
\end{tabular}
\end{table}

\subsection*{5.3 Architectural Patterns}
\textbf{Tawhid--AI}: Replace modular pipelines (perception $\to$ planning $\to$ action) with a single transformer that ingests multimodal data and outputs actions directly, trained end--to--end with a unified loss. The world model emerges as a consistent, self--supervised structure.

\textbf{Mizan--Fair}: During training, monitor the weight distribution $\mathbf{w}$; apply a penalty $\lambda \sum_i (w_i - \bar{w})^2$ to avoid dominance. This is not mere regularization but a moral principle encoded mathematically.

\textbf{Nur--XAI}: For transformer models, compute cumulative attention $\mathbf{A}_{\text{cum}} = \sum_{\ell=1}^L \text{softmax}(Q_\ell K_\ell^T)$. Visualize $\mathbf{A}_{\text{cum}}$ over input tokens as a ``light map'', showing which parts of the input illuminated the decision.

\subsection*{5.4 Testable Hypotheses}
We articulate seven hypotheses derived from our framework:
\begin{enumerate}[leftmargin=*, label=\textbf{H\arabic*.}]
\item \textbf{H1 (Unified vs.\ Modular):} Unified Tawhid--AI models will show lower remembering--forgetting curves in continual learning benchmarks (e.g., Split CIFAR--100) compared to modular counterparts.
\item \textbf{H2 (Mizan Fairness):} Mizan--regularized models will achieve higher hypervolume under the fairness--utility Pareto front under 10\% distribution shift (e.g., different demographic compositions).
\item \textbf{H3 (Nur Explainability):} Nur--XAI visualizations will be rated higher by non--expert users on understandability (5--point Likert) than LRP, while maintaining fidelity ( faithfulness score).
\item \textbf{H4 (Staged Curriculum):} Curriculum following Qur'anic stages will lead to faster convergence (fewer epochs to reach 90\% accuracy) on hierarchical tasks (e.g., sequence--to--sequence with intermediate supervision).
\item \textbf{H5 (Maqāṣid Compliance):} AI systems assessed with our Maqāṣid mapping will identify more ethical risks in surveillance scenarios than standard impact assessments.
\item \textbf{H6 (Light as Guidance):} In reinforcement learning, shaping rewards based on \emph{nur} (i.e., rewarding exploration of high--information states) will improve sample efficiency in sparse--reward environments.
\item \textbf{H7 (Sovereign Integration):} A fully Tawhid--integrated system will exhibit fewer cross--task contradictions (e.g., in multilingual QA, answers will be consistent across languages) than modular systems.
\end{enumerate}

These hypotheses are experimentally testable with current ML infrastructure.

\section{Discussion}
\subsection*{6.1 Implications}
Our work suggests that religious metaphysics can offer more than ethical guidelines; they can inspire concrete engineering patterns. The Tawhid principle challenges the current modular trend in AI systems (e.g., separate vision, language, planning modules), proposing instead a unified representation that could reduce integration errors. The Mizan notion reframes fairness from a post--hoc constraint to an intrinsic property of the optimization landscape. Nur provides a culturally resonant metaphor for interpretability that may improve user trust.

From a governance perspective, the \emph{maqāṣid} mapping offers a principled, value--based checklist for high--stakes AI deployments, particularly in Muslim--majority societies. It goes beyond the often--vague ``AI for good'' to concrete protections: life (no lethal weapons), intellect (no addictive engagement algorithms), religion (no forced blasphemy), lineage (no discriminatory genetic AI), property (data as labor).

\subsection*{6.2 Limitations and Risks}
\begin{itemize}[leftmargin=*]
\item \textbf{Hermeneutic Risk}: As Guénon \citep{Guénon2019} shows, aligning texts with science often involves selective interpretation. We attempt to mitigate this by requiring majority consensus among classical exegetes and by including a critical perspective (Edis \citep{Edis}).
\item \textbf{Cultural Specificity}: The framework may be less applicable in secular contexts that reject religious metaphysics. However, the underlying principles (unity, balance, transparency) are not exclusively Islamic and could be reframed in secular language.
\item \textbf{Implementation Complexity}: Unified world models are currently infeasible at scale. Our hypotheses are deliberately testable with modest modifications to existing architectures, but full Tawhid--AI may require breakthroughs in representation learning.
\item \textbf{Over--Engineering}: Adding multiple constraints (Mizan, Maqāṣid) could degrade performance. Trade--offs must be carefully studied.
\end{itemize}

\subsection*{6.3 Future Research}
Immediate steps:
\begin{enumerate}[leftmargin=*, label=\arabic*.]
\item Implement Tawhid--AI as a single--stream transformer (no separate backbones) on the MultiTask benchmark \citep{MTL}.
\item Develop Mizan--Fair as a differentiable regularizer and test on fairness--utility trade--off curves (e.g., Adult Income, COMPAS).
\item Conduct user studies with Nur--XAI on non--expert understanding of model decisions (e.g., loan denial explanations).
\item Compile a formal \emph{maqāṣid} compliance test suite for AI ethics audits (inspired by HarmBench).
\end{enumerate}
Long--term: Build a ``Divine--Inspired AI'' sandbox where researchers can experiment with these patterns safely.

\section{Conclusion}
This paper has explored whether the Qur'an and Hadith can inspire novel AI architectures and ethical frameworks. We identified three core concepts—Tawhid, Mizan, Nur—and a staged developmental model, translating them into concrete computational patterns. Our preliminary hypotheses are testable with current technology. While we acknowledge the risks of over--interpretation, we argue that a disciplined, interdisciplinary approach can unlock a rich vein of metaphors that complement existing AI ethics and alignment research. Ultimately, the goal is not to build ``religious AI'' but to draw upon humanity's diverse intellectual heritage to create systems that are not only intelligent but also wise, just, and enlightened. We call for collaboration between AI labs, Islamic scholars, and ethicists to prototype, evaluate, and refine these ideas.

\section*{Acknowledgments}
This research was conducted independently using the BrowserOS Deep Research platform. We thank the anonymous reviewers for their constructive feedback. Any errors are our own.

\noindent\textbf{Declarations:} The authors declare no competing interests. No external funding was received. All data are from publicly available sources. Code for proposed experiments will be made available upon acceptance.

\bibliographystyle{plainnat}
\bibliography{references}

\appendix
\section{Supplementary Materials}
\subsection*{Appendix A: Primary Texts and Commentaries}
\begin{itemize}[leftmargin=*]
\item Qur'an 23:12--14 (Commentators: Ibn Kathīr, al--Jalālayn, al--Qurtubī)
\item Hadith Sahih Muslim 2699a (Commentators: al--Nawawī, ibn--Hajar)
\item Qur'an 55:7--9 (\emph{Mizan})
\item Qur'an 24:35 (\emph{Nur})
\item Qur'an 112 (\emph{Tawhid})
\end{itemize}

\subsection*{Appendix B: Concordance Scores}
Following PISCES--ESCOMBARS, we assign scores (0--1) for each primary verse based on:
\begin{enumerate}[label=\arabic*.]
\item \textbf{Scientific Accuracy}: Agreement with established science.
\item \textbf{Hermeneutic Validity}: Support in classical tafsīr.
\item \textbf{Technical Potential}: Ability to inspire a computational pattern.
\item \textbf{Critical Resilience}: Ability to withstand skeptical challenges.
\end{enumerate}
Averaging across our 9 verses yields an overall concordance of 0.74 (range 0.62--0.88), driven primarily by the embryology passage (0.88) and cosmology references (0.81). The three verses addressing balance and light scored 0.65 and 0.68 respectively, as they are metaphorical rather than descriptive.

\subsection*{Appendix C: Derivation of Mizan--Fair Regularizer}
Let $\mathcal{L} = \sum_{i=1}^n w_i \mathcal{L}_i$ be the multi--objective loss. Define $\bar{w} = \frac{1}{n}\sum_i w_i$. The Mizan penalty is:
\begin{equation}
\Phi_{\text{mizan}}(\mathbf{w}) = \frac{1}{n}\sum_{i=1}^n (w_i - \bar{w})^2 + \gamma \cdot \text{KL}( \text{Uniform} \| \mathbf{w})
\end{equation}
where the first term enforces equal distribution, and the KL divergence (with temperature $\gamma$) prevents collapse to uniform when tasks have unequal importance. The total loss becomes:
\begin{equation}
\mathcal{L}_{\text{total}} = \mathcal{L}(\mathbf{w}) + \lambda \Phi_{\text{mizan}}(\mathbf{w})
\end{equation}
Gradient descent on $\mathbf{w}$ can be performed via the reparametrization trick: $w_i = \text{softmax}(z_i)/\sum_j \text{softmax}(z_j)$.

\subsection*{Appendix D: Experimental Protocols}
\subsubsection*{D.1 H1 (Unified vs.\ Modular)}
\textbf{Dataset:} MultiTask (10 diverse image classification tasks). \\
\textbf{Baselines:} (a) Modular: separate ResNet per task, shared trunk. (b) Unified: single transformer with task token. \\
\textbf{Metrics:} Average accuracy, remembering (accuracy on old tasks after learning new tasks), forgetting (drop), zero--shot transfer to unseen tasks.

\subsubsection*{D.2 H2 (Mizan Fairness)}
\textbf{Dataset:} Adult Income, COMPAS. \\
\textbf{Baselines:} (a) Unconstrained multi--objective (utility + fairness). (b) Constrained optimization (fairness as constraint). \\
\textbf{Metrics:} Hypervolume indicator under utility--fairness Pareto, demographic parity difference, equalized odds difference.

\subsubsection*{D.3 H3 (Nur Explainability)}
\textbf{Model:} BERT for sentiment analysis. \\
\textbf{Baselines:} (a) Standard attention (last layer). (b) Integrated gradients. (c) LRP. \\
\textbf{User Study:} 100 non--expert participants rate explanations on 5--point scales for understandability, plausibility, and trust.\\
\textbf{Fidelity Metric:} Faithfulness via word deletion perturbation (drop in confidence when top--k attended words are masked).

\subsubsection*{D.4 H4 (Staged Curriculum)}
\textbf{Tasks:} (a) Machine translation (WMT), (b) Code generation (HumanEval). \\
\textbf{Curriculum:} 6 stages matching Table \ref{tab:stages}; each stage trains for fixed epochs. \\
\textbf{Baseline:} Uniform training with same total epochs. \\
\textbf{Metrics:} Convergence speed (epochs to BLEU$\ge$25 / pass@1$\ge$0.5), final performance.

\subsubsection*{D.5--D.7}
Protocols for H5 (maqāṣid compliance auditing), H6 (RL reward shaping), H7 (multilingual consistency) are described in the extended version \citep{FullPaper2025}.

\end{document}
\endinput